% mazzi: 24
\chapter{L'apparato genitale femminile}

% ---------------------------------------------------------------------------
\section{Generalità}

\textbf{Apparato genitale femminile}: insieme di organi e strutture che permettono la
riproduzione e l'accoppiamento.

\rubrica{Costituzione}
\begin{itemize}
  \item \textbf{organi genitali}: \textbf{ovaie};
  \item \textbf{vie genitali}: \textbf{tube uterine}, \textbf{utero}, \textbf{vagina};
  \item \textbf{genitali esterni}: \textbf{vulva} (grandi e piccole labbra, più ghiandole).
\end{itemize}

\rubrica{Confronto con l'apparato genitale maschile}
\begin{itemize}
  \item genitali interni:
  \begin{itemize}
    \item maschili: testicolo, epididimo, dotto deferente, prostata, vescichetta seminale,
      ghiandola bulbouretrale;
    \item femminili: ovaio, utero, tuba uterina, vagina (parte superiore);
  \end{itemize}
  \item genitali esterni:
  \begin{itemize}
    \item maschili: pene con uretra, scroto con gli involucri del testicolo;
    \item femminili: vagina (soltanto il vestibolo), grandi/piccole labbra, monte di Venere,
      ghiandole vestibolari maggiore/minore, clitoride.
  \end{itemize}
\end{itemize}

% ---------------------------------------------------------------------------
\section{Utero}

\textbf{Utero}: organo cavo, impari, mediale, nella cavità pelvica, parzialmente coperto da
peritoneo; davanti al retto e postero-superiormente alla vescica.
\begin{itemize}
  \item forma di pera capovolta; dimensioni: $7{,}5 \times 5 \times 2$ cm; peso: $\sim$90 g;
    solitamente antiverso e antiflesso;
  \item strutturato in: un \textbf{fondo} (sopra lo sbocco delle tube), un \textbf{corpo}, un
    \textbf{collo} (porzione intra- e sopravaginale), 2 facce (anteriore e posteriore), 2 margini
    laterali, 1 cavità uterina in comunicazione con il canale cervicale.
\end{itemize}

Il \textbf{fondo} si trova superiormente allo sbocco delle tube uterine. Il \textbf{collo}
termina sporgendo dentro la vagina (porzione intravaginale del collo, detta ``a muso di
tinca''), dove termina con l'\textbf{orifizio esterno}. Il \textbf{canale cervicale} presenta un
orifizio interno e uno esterno, col quale si apre in vagina.

\rubrica{Antiversione e antiflessione}
Due angoli aperti anteriormente, tra:
\begin{itemize}
  \item corpo e collo (\textbf{antiflessione});
  \item collo (utero) e vagina (\textbf{antiversione});
  \item il collo dell'utero entra in vagina guardando la sua parete posteriore, presentando la
    sporgenza a muso di tinca con l'orifizio esterno;
  \item la vagina forma dei recessi, detti \textbf{fornici}, attorno al collo uterino (anteriore,
    posteriore, laterali).
\end{itemize}

\rubrica{Rapporti peritoneali}
Utero parzialmente rivestito dal peritoneo.
\begin{itemize}
  \item le pieghe del peritoneo formano in avanti il \textbf{cavo utero-vescicale}, in addietro il
    \textbf{cavo utero-rettale} (\textit{cavo del Douglas}: estensione della cavità peritoneale tra
    retto e parete posteriore dell'utero);
  \item il peritoneo che copre l'utero si continua lateralmente fino alle pareti laterali della
    piccola pelvi, formando il \textbf{legamento largo dell'utero}: legamento peritoneale,
    disposto frontalmente, che collega il margine laterale dell'utero alla parete pelvica; presenta
    2 facce, anteriore e posteriore.
\end{itemize}

% ---------------------------------------------------------------------------
\section{Parete uterina}

\begin{itemize}
  \item \textbf{endometrio}: tonaca mucosa con epitelio cilindrico semplice e ghiandole; varia
    molto durante le fasi del ciclo uterino;
  \item \textbf{miometrio}: tonaca muscolare liscia in più strati, molto spessa e abbondantemente
    vascolarizzata;
  \item \textbf{perimetrio}: in parte dal peritoneo che lo ricopre, in parte dal connettivo che lo
    ricopre.
\end{itemize}

\rubrica{Endometrio}
Tonaca mucosa, con epitelio batiprismatico semplice, cellule ciliate e cellule secernenti
alternate; verso la vagina, e nella vagina stessa, l'epitelio diventa pavimentoso composto.
\begin{itemize}
  \item nella lamina basale si trovano le numerose \textbf{ghiandole uterine};
  \item nel corpo dell'utero (non nel collo) l'endometrio subisce continue modificazioni cicliche:
    il \textbf{ciclo mestruale}.
\end{itemize}

\rubrica{Miometrio}
Spesso strato muscolare, in 3 strati (più distinguibili nel corpo che nel collo dell'utero):
\begin{itemize}
  \item \textbf{sottomucoso}: sottile, fibrocellule lisce a decorso longitudinale e obliquo;
  \item \textbf{vascolare}: molto spesso, decorso circolare, vi decorrono i vasi principali
    dell'utero; la contrazione dà emostasi;
  \item \textbf{sottosieroso}: sottile, decorso longitudinale.
\end{itemize}

% ---------------------------------------------------------------------------
\section{Il ciclo mestruale}

Modificazione dell'endometrio del corpo dell'utero, per effetto degli ormoni ovarici
\textbf{estrogeni} e \textbf{progesterone}; dura $\sim$28 giorni, in 4 fasi successive:
\textbf{rigenerativa}, \textbf{proliferativa}, \textbf{secretiva}, \textbf{desquamativa}.

\rubrica{Fase rigenerativa e fase proliferativa}
Nell'ovaio matura un follicolo ooforo, che rilascia estrogeni.
\begin{itemize}
  \item \textbf{rigenerativa}: si rigenera l'epitelio dell'endometrio, a partire dai fondi delle
    ghiandole, dalle quali l'epitelio prolifera per coprire lo stroma rimasto nudo;
  \item \textbf{proliferativa}: l'endometrio si fa più spesso, appaiono le cellule ciliate, le
    ghiandole sono più lunghe.
\end{itemize}

\rubrica{Fase secretiva}
Nell'ovaio è avvenuta l'ovulazione e il corpo luteo è attivo: rilascia progesterone, quindi c'è
l'effetto combinato di estrogeni e progesterone. I vasi sanguiferi della lamina propria si fanno
molto dilatati, le ghiandole si dilatano in profondità accumulando un secreto mucoso.

\rubrica{Fase desquamativa}
Non è avvenuto l'impianto dell'embrione.
\begin{itemize}
  \item il corpo luteo si degrada e viene a mancare la sintesi di progesterone;
  \item si desquamano e degradano le ghiandole della lamina basale della mucosa uterina;
  \item \textbf{flusso mestruale}: si formano trombi nei vasi della lamina propria, quindi ridotta
    irrorazione sanguigna, necrosi e desquamazione dell'endometrio.
\end{itemize}

% ---------------------------------------------------------------------------
\section{Apparato di sospensione dell'utero}

Il \textbf{legamento ovarico} è un legamento fibroso che connette la parte inferiore-mediale
dell'ovaio alla superficie laterale dell'utero.

Il collo è la parte meno mobile dell'utero.

\rubrica{Legamenti del collo dell'utero}
\begin{enumerate}
  \item ad andamento trasversale: \textbf{legamento cardinale} (o di \textit{Mackenrodt}), dal
    collo alle pareti laterali della pelvi;
  \item ad andamento sagittale:
  \begin{itemize}
    \item \textbf{uterosacrali}: dalla parete posteriore, si dirigono superiormente e
      posteriormente al promontorio sacro-vertebrale, circondando il retto;
    \item \textbf{uteropubici}: dalla parete anteriore al pube, passando lateralmente alla
      vescica.
  \end{itemize}
\end{enumerate}

L'utero è vascolarizzato dall'\textbf{arteria uterina} (ramo dell'arteria iliaca interna, o
ipogastrica), che irrora anche uretere e vescica urinaria.

% ---------------------------------------------------------------------------
\section{Tube uterine}

\textbf{Tube uterine} (o \textit{di Falloppio}): lunghezza 12--18 cm ($\sim$14--15 cm), spessore
$\sim$3 mm.
\begin{itemize}
  \item più ampie lateralmente, si assottigliano verso l'utero;
  \item vanno dal polo superiore dell'ovaio all'angolo tubarico dell'utero;
  \item parti: \textbf{infundibolo} (con fimbrie, estremità laterale), \textbf{ampolla},
    \textbf{istmo}, \textbf{parte interstiziale} (attraversa la spessa parete dell'utero);
  \item sono peritoneali; rivestite dal peritoneo, sono collegate alla cavità peritoneale da
    foglietti peritoneali detti \textbf{mesosalpingi}; la piega del legamento largo dell'utero sulle
    tube forma la \textbf{borsa ovarica}, ove cade l'ovulo maturo.
\end{itemize}

\rubrica{Conformazione interna}
\begin{itemize}
  \item numerose pieghe (sezione stellata) che si anastomizzano tra loro;
  \item epitelio cilindrico, con cellule ciliate (trasporto dell'ovulo) e altre più basse
    secernenti muco;
  \item parete muscolare: contrazioni per la propulsione dell'ovulo.
\end{itemize}
Gli spermatozoi risalgono fino alle tube uterine: l'uovo è fecondato nelle tube uterine
(nell'ampolla) e raggiunge l'utero solo dopo i primi stadi di divisione cellulare.

% ---------------------------------------------------------------------------
\section{Ovaie}

Le \textbf{ovaie} sono intraperitoneali, nella cavità pelvica.
\begin{itemize}
  \item \textbf{margine anteriore}: \textbf{meso ovarico}, piega peritoneale che proviene dal
    legamento largo dell'utero, veicola vasi e nervi all'interno dell'ovaio;
  \item \textbf{polo superiore}: \textbf{legamento sospensore dell'ovaio} (veicola vasi e nervi) e
    \textbf{legamento tubo-ovarico} (lega l'ovaio alla tuba);
  \item \textbf{polo inferiore}: \textbf{legamento proprio dell'ovaio} (o utero-ovarico), va
    all'angolo tubarico dell'utero;
  \item circondate dalla \textbf{borsa ovarica}: tasca di origine peritoneale.
\end{itemize}

\rubrica{Struttura}
Il parenchima delle ovaie si distingue in:
\begin{itemize}
  \item \textbf{corticale}: più esterna, contenente i gameti in formazione; costituita da
    connettivo (che in periferia forma una capsula fibrosa, la \textbf{tonaca albuginea}) ed
    \textbf{epitelio germinativo}; qui si trovano i \textbf{follicoli}, complesse strutture in
    vari stadi di maturazione, formate da diversi strati di cellule contenenti l'oocita in
    differenziamento;
  \item \textbf{midollare}: più interna, contenente vasi e nervi; costituita da connettivo fibroso
    e cellule muscolari lisce.
\end{itemize}

\rubrica{Funzioni e irrorazione}
\begin{itemize}
  \item funzioni: \textbf{riproduttiva} ed \textbf{endocrina};
  \item vascolarizzazione: arteria ovarica e arteria uterina;
  \item innervazione: plesso utero-ovarico e fibre nervose dal plesso pelvico.
\end{itemize}

% ---------------------------------------------------------------------------
\section{Vagina}

\textbf{Vagina}: parte elastica e muscolare del tratto genitale femminile; canale
fibro-muscolare molto elastico che serve da supporto al collo dell'utero e all'uretra. Va dalla
vulva alla cervice; lunghezza $\sim$7,5--9 cm.

\rubrica{Struttura}
Il canale vaginale è rivestito da una membrana mucosa costituita superficialmente da un
\textbf{epitelio pavimentoso pluristratificato}: numerosi strati di cellule, che nello strato più
esterno assumono una forma appiattita, più sviluppata in lunghezza. 3 strati di tessuto:
\textbf{epitelio squamoso}, \textbf{muscolare liscio}, \textbf{avventizia}.

% ---------------------------------------------------------------------------
\section{Genitali esterni (vulva)}

Struttura comprendente: monte del pube, grandi labbra, piccole labbra, prepuzio e glande del
clitoride, orifizio uretrale, vestibolo, orifizio della vagina (protetto dall'\textbf{imene}).

% ---------------------------------------------------------------------------
\section{Confronto tra spermatogenesi e ovogenesi}

\begin{itemize}
  \item \textbf{spermatogenesi}: processo continuo; 1 spermatocita origina 4 spermatozoi;
  \item \textbf{ovogenesi}: processo ciclico; 1 ovocita origina 1 uovo.
\end{itemize}
Rilevanza clinica: nella donna, agenti quali le radiazioni ionizzanti inducono sterilità
permanente.

% ---------------------------------------------------------------------------
\section{Mammella}

\textbf{Mammella}: organo ghiandolare che, nelle femmine di mammifero, secerne il latte;
caratteristica sessuale secondaria nella donna, con anche una valenza simbolica.

\rubrica{Struttura}
\begin{itemize}
  \item \textbf{componente ghiandolare}: 15--20 \textbf{lobi}, ciascuno con uno sbocco verso il
    capezzolo attraverso un \textbf{dotto galattoforo};
  \item \textbf{componente adiposa}: in cui sono inserite e immerse le strutture ghiandolari;
  \item \textbf{componente fibrosa} di sostegno: genera suddivisioni tra le diverse appendici
    ghiandolari.
\end{itemize}
